% gramatica_beamer.tex
% Slide da gramática (problema Smallest, versão simplificada).
% Inclua no main.tex com % gramatica_beamer.tex
% Slide da gramática (problema Smallest, versão simplificada).
% Inclua no main.tex com % gramatica_beamer.tex
% Slide da gramática (problema Smallest, versão simplificada).
% Inclua no main.tex com % gramatica_beamer.tex
% Slide da gramática (problema Smallest, versão simplificada).
% Inclua no main.tex com \input{gramatica_beamer} ou \include{gramatica_beamer}.
% Não precisa de pacotes extras: só do \usepackage[T1]{fontenc} (já presente
% no main), para que < e > saiam corretos em fonte monoespaçada.

% Token especial (NEWLINE, INDENT) em negrito sem serifa
\providecommand{\tok}[1]{\mbox{\textsf{\textbf{#1}}}}
% Barra que separa as alternativas da BNF
\providecommand{\ou}{$\mid$\ }

\begin{frame}{\textit{Grammar-guided} GP}
  \begin{figure}
  \centering
  \begin{examples}
  \small
  \ttfamily
  \setlength{\tabcolsep}{4pt}
  \renewcommand{\arraystretch}{1.35}
  \begin{tabular}{@{}l l l@{}}
    &n = &[13, 17, 7, 24, 21, 20, 18, 28, 4, 29] \\
    <program>  & ::= & def sum(int1, int2): \tok{NEWLINE} \tok{INDENT} \\
               &     & int3 = int() \tok{NEWLINE} \\
               &     & <stmt> return <int> \\
    <stmt>     & ::= & <stmt><stmt> \\
               & \ou & <int\_var> = <int> \tok{NEWLINE} \\
    <int\_var> & ::= & int1 \ou int2 \ou int3 \\
    <int\_op>  & ::= & + \ou - \ou * \ou / \ou \% \\
    <int>      & ::= & <int\_var> \\
               & \ou & <int> <int\_op> <int> \\
  \end{tabular}
  \end{examples}
  \caption{Exemplo de uma derivação de programa GE. Adaptado de \cite{survey}.}
  \end{figure}
\end{frame}
 ou % gramatica_beamer.tex
% Slide da gramática (problema Smallest, versão simplificada).
% Inclua no main.tex com \input{gramatica_beamer} ou \include{gramatica_beamer}.
% Não precisa de pacotes extras: só do \usepackage[T1]{fontenc} (já presente
% no main), para que < e > saiam corretos em fonte monoespaçada.

% Token especial (NEWLINE, INDENT) em negrito sem serifa
\providecommand{\tok}[1]{\mbox{\textsf{\textbf{#1}}}}
% Barra que separa as alternativas da BNF
\providecommand{\ou}{$\mid$\ }

\begin{frame}{\textit{Grammar-guided} GP}
  \begin{figure}
  \centering
  \begin{examples}
  \small
  \ttfamily
  \setlength{\tabcolsep}{4pt}
  \renewcommand{\arraystretch}{1.35}
  \begin{tabular}{@{}l l l@{}}
    &n = &[13, 17, 7, 24, 21, 20, 18, 28, 4, 29] \\
    <program>  & ::= & def sum(int1, int2): \tok{NEWLINE} \tok{INDENT} \\
               &     & int3 = int() \tok{NEWLINE} \\
               &     & <stmt> return <int> \\
    <stmt>     & ::= & <stmt><stmt> \\
               & \ou & <int\_var> = <int> \tok{NEWLINE} \\
    <int\_var> & ::= & int1 \ou int2 \ou int3 \\
    <int\_op>  & ::= & + \ou - \ou * \ou / \ou \% \\
    <int>      & ::= & <int\_var> \\
               & \ou & <int> <int\_op> <int> \\
  \end{tabular}
  \end{examples}
  \caption{Exemplo de uma derivação de programa GE. Adaptado de \cite{survey}.}
  \end{figure}
\end{frame}
.
% Não precisa de pacotes extras: só do \usepackage[T1]{fontenc} (já presente
% no main), para que < e > saiam corretos em fonte monoespaçada.

% Token especial (NEWLINE, INDENT) em negrito sem serifa
\providecommand{\tok}[1]{\mbox{\textsf{\textbf{#1}}}}
% Barra que separa as alternativas da BNF
\providecommand{\ou}{$\mid$\ }

\begin{frame}{\textit{Grammar-guided} GP}
  \begin{figure}
  \centering
  \begin{examples}
  \small
  \ttfamily
  \setlength{\tabcolsep}{4pt}
  \renewcommand{\arraystretch}{1.35}
  \begin{tabular}{@{}l l l@{}}
    &n = &[13, 17, 7, 24, 21, 20, 18, 28, 4, 29] \\
    <program>  & ::= & def sum(int1, int2): \tok{NEWLINE} \tok{INDENT} \\
               &     & int3 = int() \tok{NEWLINE} \\
               &     & <stmt> return <int> \\
    <stmt>     & ::= & <stmt><stmt> \\
               & \ou & <int\_var> = <int> \tok{NEWLINE} \\
    <int\_var> & ::= & int1 \ou int2 \ou int3 \\
    <int\_op>  & ::= & + \ou - \ou * \ou / \ou \% \\
    <int>      & ::= & <int\_var> \\
               & \ou & <int> <int\_op> <int> \\
  \end{tabular}
  \end{examples}
  \caption{Exemplo de uma derivação de programa GE. Adaptado de \cite{survey}.}
  \end{figure}
\end{frame}
 ou % gramatica_beamer.tex
% Slide da gramática (problema Smallest, versão simplificada).
% Inclua no main.tex com % gramatica_beamer.tex
% Slide da gramática (problema Smallest, versão simplificada).
% Inclua no main.tex com \input{gramatica_beamer} ou \include{gramatica_beamer}.
% Não precisa de pacotes extras: só do \usepackage[T1]{fontenc} (já presente
% no main), para que < e > saiam corretos em fonte monoespaçada.

% Token especial (NEWLINE, INDENT) em negrito sem serifa
\providecommand{\tok}[1]{\mbox{\textsf{\textbf{#1}}}}
% Barra que separa as alternativas da BNF
\providecommand{\ou}{$\mid$\ }

\begin{frame}{\textit{Grammar-guided} GP}
  \begin{figure}
  \centering
  \begin{examples}
  \small
  \ttfamily
  \setlength{\tabcolsep}{4pt}
  \renewcommand{\arraystretch}{1.35}
  \begin{tabular}{@{}l l l@{}}
    &n = &[13, 17, 7, 24, 21, 20, 18, 28, 4, 29] \\
    <program>  & ::= & def sum(int1, int2): \tok{NEWLINE} \tok{INDENT} \\
               &     & int3 = int() \tok{NEWLINE} \\
               &     & <stmt> return <int> \\
    <stmt>     & ::= & <stmt><stmt> \\
               & \ou & <int\_var> = <int> \tok{NEWLINE} \\
    <int\_var> & ::= & int1 \ou int2 \ou int3 \\
    <int\_op>  & ::= & + \ou - \ou * \ou / \ou \% \\
    <int>      & ::= & <int\_var> \\
               & \ou & <int> <int\_op> <int> \\
  \end{tabular}
  \end{examples}
  \caption{Exemplo de uma derivação de programa GE. Adaptado de \cite{survey}.}
  \end{figure}
\end{frame}
 ou % gramatica_beamer.tex
% Slide da gramática (problema Smallest, versão simplificada).
% Inclua no main.tex com \input{gramatica_beamer} ou \include{gramatica_beamer}.
% Não precisa de pacotes extras: só do \usepackage[T1]{fontenc} (já presente
% no main), para que < e > saiam corretos em fonte monoespaçada.

% Token especial (NEWLINE, INDENT) em negrito sem serifa
\providecommand{\tok}[1]{\mbox{\textsf{\textbf{#1}}}}
% Barra que separa as alternativas da BNF
\providecommand{\ou}{$\mid$\ }

\begin{frame}{\textit{Grammar-guided} GP}
  \begin{figure}
  \centering
  \begin{examples}
  \small
  \ttfamily
  \setlength{\tabcolsep}{4pt}
  \renewcommand{\arraystretch}{1.35}
  \begin{tabular}{@{}l l l@{}}
    &n = &[13, 17, 7, 24, 21, 20, 18, 28, 4, 29] \\
    <program>  & ::= & def sum(int1, int2): \tok{NEWLINE} \tok{INDENT} \\
               &     & int3 = int() \tok{NEWLINE} \\
               &     & <stmt> return <int> \\
    <stmt>     & ::= & <stmt><stmt> \\
               & \ou & <int\_var> = <int> \tok{NEWLINE} \\
    <int\_var> & ::= & int1 \ou int2 \ou int3 \\
    <int\_op>  & ::= & + \ou - \ou * \ou / \ou \% \\
    <int>      & ::= & <int\_var> \\
               & \ou & <int> <int\_op> <int> \\
  \end{tabular}
  \end{examples}
  \caption{Exemplo de uma derivação de programa GE. Adaptado de \cite{survey}.}
  \end{figure}
\end{frame}
.
% Não precisa de pacotes extras: só do \usepackage[T1]{fontenc} (já presente
% no main), para que < e > saiam corretos em fonte monoespaçada.

% Token especial (NEWLINE, INDENT) em negrito sem serifa
\providecommand{\tok}[1]{\mbox{\textsf{\textbf{#1}}}}
% Barra que separa as alternativas da BNF
\providecommand{\ou}{$\mid$\ }

\begin{frame}{\textit{Grammar-guided} GP}
  \begin{figure}
  \centering
  \begin{examples}
  \small
  \ttfamily
  \setlength{\tabcolsep}{4pt}
  \renewcommand{\arraystretch}{1.35}
  \begin{tabular}{@{}l l l@{}}
    &n = &[13, 17, 7, 24, 21, 20, 18, 28, 4, 29] \\
    <program>  & ::= & def sum(int1, int2): \tok{NEWLINE} \tok{INDENT} \\
               &     & int3 = int() \tok{NEWLINE} \\
               &     & <stmt> return <int> \\
    <stmt>     & ::= & <stmt><stmt> \\
               & \ou & <int\_var> = <int> \tok{NEWLINE} \\
    <int\_var> & ::= & int1 \ou int2 \ou int3 \\
    <int\_op>  & ::= & + \ou - \ou * \ou / \ou \% \\
    <int>      & ::= & <int\_var> \\
               & \ou & <int> <int\_op> <int> \\
  \end{tabular}
  \end{examples}
  \caption{Exemplo de uma derivação de programa GE. Adaptado de \cite{survey}.}
  \end{figure}
\end{frame}
.
% Não precisa de pacotes extras: só do \usepackage[T1]{fontenc} (já presente
% no main), para que < e > saiam corretos em fonte monoespaçada.

% Token especial (NEWLINE, INDENT) em negrito sem serifa
\providecommand{\tok}[1]{\mbox{\textsf{\textbf{#1}}}}
% Barra que separa as alternativas da BNF
\providecommand{\ou}{$\mid$\ }

\begin{frame}{\textit{Grammar-guided} GP}
  \begin{figure}
  \centering
  \begin{examples}
  \small
  \ttfamily
  \setlength{\tabcolsep}{4pt}
  \renewcommand{\arraystretch}{1.35}
  \begin{tabular}{@{}l l l@{}}
    &n = &[13, 17, 7, 24, 21, 20, 18, 28, 4, 29] \\
    <program>  & ::= & def sum(int1, int2): \tok{NEWLINE} \tok{INDENT} \\
               &     & int3 = int() \tok{NEWLINE} \\
               &     & <stmt> return <int> \\
    <stmt>     & ::= & <stmt><stmt> \\
               & \ou & <int\_var> = <int> \tok{NEWLINE} \\
    <int\_var> & ::= & int1 \ou int2 \ou int3 \\
    <int\_op>  & ::= & + \ou - \ou * \ou / \ou \% \\
    <int>      & ::= & <int\_var> \\
               & \ou & <int> <int\_op> <int> \\
  \end{tabular}
  \end{examples}
  \caption{Exemplo de uma derivação de programa GE. Adaptado de \cite{survey}.}
  \end{figure}
\end{frame}
 ou % gramatica_beamer.tex
% Slide da gramática (problema Smallest, versão simplificada).
% Inclua no main.tex com % gramatica_beamer.tex
% Slide da gramática (problema Smallest, versão simplificada).
% Inclua no main.tex com % gramatica_beamer.tex
% Slide da gramática (problema Smallest, versão simplificada).
% Inclua no main.tex com \input{gramatica_beamer} ou \include{gramatica_beamer}.
% Não precisa de pacotes extras: só do \usepackage[T1]{fontenc} (já presente
% no main), para que < e > saiam corretos em fonte monoespaçada.

% Token especial (NEWLINE, INDENT) em negrito sem serifa
\providecommand{\tok}[1]{\mbox{\textsf{\textbf{#1}}}}
% Barra que separa as alternativas da BNF
\providecommand{\ou}{$\mid$\ }

\begin{frame}{\textit{Grammar-guided} GP}
  \begin{figure}
  \centering
  \begin{examples}
  \small
  \ttfamily
  \setlength{\tabcolsep}{4pt}
  \renewcommand{\arraystretch}{1.35}
  \begin{tabular}{@{}l l l@{}}
    &n = &[13, 17, 7, 24, 21, 20, 18, 28, 4, 29] \\
    <program>  & ::= & def sum(int1, int2): \tok{NEWLINE} \tok{INDENT} \\
               &     & int3 = int() \tok{NEWLINE} \\
               &     & <stmt> return <int> \\
    <stmt>     & ::= & <stmt><stmt> \\
               & \ou & <int\_var> = <int> \tok{NEWLINE} \\
    <int\_var> & ::= & int1 \ou int2 \ou int3 \\
    <int\_op>  & ::= & + \ou - \ou * \ou / \ou \% \\
    <int>      & ::= & <int\_var> \\
               & \ou & <int> <int\_op> <int> \\
  \end{tabular}
  \end{examples}
  \caption{Exemplo de uma derivação de programa GE. Adaptado de \cite{survey}.}
  \end{figure}
\end{frame}
 ou % gramatica_beamer.tex
% Slide da gramática (problema Smallest, versão simplificada).
% Inclua no main.tex com \input{gramatica_beamer} ou \include{gramatica_beamer}.
% Não precisa de pacotes extras: só do \usepackage[T1]{fontenc} (já presente
% no main), para que < e > saiam corretos em fonte monoespaçada.

% Token especial (NEWLINE, INDENT) em negrito sem serifa
\providecommand{\tok}[1]{\mbox{\textsf{\textbf{#1}}}}
% Barra que separa as alternativas da BNF
\providecommand{\ou}{$\mid$\ }

\begin{frame}{\textit{Grammar-guided} GP}
  \begin{figure}
  \centering
  \begin{examples}
  \small
  \ttfamily
  \setlength{\tabcolsep}{4pt}
  \renewcommand{\arraystretch}{1.35}
  \begin{tabular}{@{}l l l@{}}
    &n = &[13, 17, 7, 24, 21, 20, 18, 28, 4, 29] \\
    <program>  & ::= & def sum(int1, int2): \tok{NEWLINE} \tok{INDENT} \\
               &     & int3 = int() \tok{NEWLINE} \\
               &     & <stmt> return <int> \\
    <stmt>     & ::= & <stmt><stmt> \\
               & \ou & <int\_var> = <int> \tok{NEWLINE} \\
    <int\_var> & ::= & int1 \ou int2 \ou int3 \\
    <int\_op>  & ::= & + \ou - \ou * \ou / \ou \% \\
    <int>      & ::= & <int\_var> \\
               & \ou & <int> <int\_op> <int> \\
  \end{tabular}
  \end{examples}
  \caption{Exemplo de uma derivação de programa GE. Adaptado de \cite{survey}.}
  \end{figure}
\end{frame}
.
% Não precisa de pacotes extras: só do \usepackage[T1]{fontenc} (já presente
% no main), para que < e > saiam corretos em fonte monoespaçada.

% Token especial (NEWLINE, INDENT) em negrito sem serifa
\providecommand{\tok}[1]{\mbox{\textsf{\textbf{#1}}}}
% Barra que separa as alternativas da BNF
\providecommand{\ou}{$\mid$\ }

\begin{frame}{\textit{Grammar-guided} GP}
  \begin{figure}
  \centering
  \begin{examples}
  \small
  \ttfamily
  \setlength{\tabcolsep}{4pt}
  \renewcommand{\arraystretch}{1.35}
  \begin{tabular}{@{}l l l@{}}
    &n = &[13, 17, 7, 24, 21, 20, 18, 28, 4, 29] \\
    <program>  & ::= & def sum(int1, int2): \tok{NEWLINE} \tok{INDENT} \\
               &     & int3 = int() \tok{NEWLINE} \\
               &     & <stmt> return <int> \\
    <stmt>     & ::= & <stmt><stmt> \\
               & \ou & <int\_var> = <int> \tok{NEWLINE} \\
    <int\_var> & ::= & int1 \ou int2 \ou int3 \\
    <int\_op>  & ::= & + \ou - \ou * \ou / \ou \% \\
    <int>      & ::= & <int\_var> \\
               & \ou & <int> <int\_op> <int> \\
  \end{tabular}
  \end{examples}
  \caption{Exemplo de uma derivação de programa GE. Adaptado de \cite{survey}.}
  \end{figure}
\end{frame}
 ou % gramatica_beamer.tex
% Slide da gramática (problema Smallest, versão simplificada).
% Inclua no main.tex com % gramatica_beamer.tex
% Slide da gramática (problema Smallest, versão simplificada).
% Inclua no main.tex com \input{gramatica_beamer} ou \include{gramatica_beamer}.
% Não precisa de pacotes extras: só do \usepackage[T1]{fontenc} (já presente
% no main), para que < e > saiam corretos em fonte monoespaçada.

% Token especial (NEWLINE, INDENT) em negrito sem serifa
\providecommand{\tok}[1]{\mbox{\textsf{\textbf{#1}}}}
% Barra que separa as alternativas da BNF
\providecommand{\ou}{$\mid$\ }

\begin{frame}{\textit{Grammar-guided} GP}
  \begin{figure}
  \centering
  \begin{examples}
  \small
  \ttfamily
  \setlength{\tabcolsep}{4pt}
  \renewcommand{\arraystretch}{1.35}
  \begin{tabular}{@{}l l l@{}}
    &n = &[13, 17, 7, 24, 21, 20, 18, 28, 4, 29] \\
    <program>  & ::= & def sum(int1, int2): \tok{NEWLINE} \tok{INDENT} \\
               &     & int3 = int() \tok{NEWLINE} \\
               &     & <stmt> return <int> \\
    <stmt>     & ::= & <stmt><stmt> \\
               & \ou & <int\_var> = <int> \tok{NEWLINE} \\
    <int\_var> & ::= & int1 \ou int2 \ou int3 \\
    <int\_op>  & ::= & + \ou - \ou * \ou / \ou \% \\
    <int>      & ::= & <int\_var> \\
               & \ou & <int> <int\_op> <int> \\
  \end{tabular}
  \end{examples}
  \caption{Exemplo de uma derivação de programa GE. Adaptado de \cite{survey}.}
  \end{figure}
\end{frame}
 ou % gramatica_beamer.tex
% Slide da gramática (problema Smallest, versão simplificada).
% Inclua no main.tex com \input{gramatica_beamer} ou \include{gramatica_beamer}.
% Não precisa de pacotes extras: só do \usepackage[T1]{fontenc} (já presente
% no main), para que < e > saiam corretos em fonte monoespaçada.

% Token especial (NEWLINE, INDENT) em negrito sem serifa
\providecommand{\tok}[1]{\mbox{\textsf{\textbf{#1}}}}
% Barra que separa as alternativas da BNF
\providecommand{\ou}{$\mid$\ }

\begin{frame}{\textit{Grammar-guided} GP}
  \begin{figure}
  \centering
  \begin{examples}
  \small
  \ttfamily
  \setlength{\tabcolsep}{4pt}
  \renewcommand{\arraystretch}{1.35}
  \begin{tabular}{@{}l l l@{}}
    &n = &[13, 17, 7, 24, 21, 20, 18, 28, 4, 29] \\
    <program>  & ::= & def sum(int1, int2): \tok{NEWLINE} \tok{INDENT} \\
               &     & int3 = int() \tok{NEWLINE} \\
               &     & <stmt> return <int> \\
    <stmt>     & ::= & <stmt><stmt> \\
               & \ou & <int\_var> = <int> \tok{NEWLINE} \\
    <int\_var> & ::= & int1 \ou int2 \ou int3 \\
    <int\_op>  & ::= & + \ou - \ou * \ou / \ou \% \\
    <int>      & ::= & <int\_var> \\
               & \ou & <int> <int\_op> <int> \\
  \end{tabular}
  \end{examples}
  \caption{Exemplo de uma derivação de programa GE. Adaptado de \cite{survey}.}
  \end{figure}
\end{frame}
.
% Não precisa de pacotes extras: só do \usepackage[T1]{fontenc} (já presente
% no main), para que < e > saiam corretos em fonte monoespaçada.

% Token especial (NEWLINE, INDENT) em negrito sem serifa
\providecommand{\tok}[1]{\mbox{\textsf{\textbf{#1}}}}
% Barra que separa as alternativas da BNF
\providecommand{\ou}{$\mid$\ }

\begin{frame}{\textit{Grammar-guided} GP}
  \begin{figure}
  \centering
  \begin{examples}
  \small
  \ttfamily
  \setlength{\tabcolsep}{4pt}
  \renewcommand{\arraystretch}{1.35}
  \begin{tabular}{@{}l l l@{}}
    &n = &[13, 17, 7, 24, 21, 20, 18, 28, 4, 29] \\
    <program>  & ::= & def sum(int1, int2): \tok{NEWLINE} \tok{INDENT} \\
               &     & int3 = int() \tok{NEWLINE} \\
               &     & <stmt> return <int> \\
    <stmt>     & ::= & <stmt><stmt> \\
               & \ou & <int\_var> = <int> \tok{NEWLINE} \\
    <int\_var> & ::= & int1 \ou int2 \ou int3 \\
    <int\_op>  & ::= & + \ou - \ou * \ou / \ou \% \\
    <int>      & ::= & <int\_var> \\
               & \ou & <int> <int\_op> <int> \\
  \end{tabular}
  \end{examples}
  \caption{Exemplo de uma derivação de programa GE. Adaptado de \cite{survey}.}
  \end{figure}
\end{frame}
.
% Não precisa de pacotes extras: só do \usepackage[T1]{fontenc} (já presente
% no main), para que < e > saiam corretos em fonte monoespaçada.

% Token especial (NEWLINE, INDENT) em negrito sem serifa
\providecommand{\tok}[1]{\mbox{\textsf{\textbf{#1}}}}
% Barra que separa as alternativas da BNF
\providecommand{\ou}{$\mid$\ }

\begin{frame}{\textit{Grammar-guided} GP}
  \begin{figure}
  \centering
  \begin{examples}
  \small
  \ttfamily
  \setlength{\tabcolsep}{4pt}
  \renewcommand{\arraystretch}{1.35}
  \begin{tabular}{@{}l l l@{}}
    &n = &[13, 17, 7, 24, 21, 20, 18, 28, 4, 29] \\
    <program>  & ::= & def sum(int1, int2): \tok{NEWLINE} \tok{INDENT} \\
               &     & int3 = int() \tok{NEWLINE} \\
               &     & <stmt> return <int> \\
    <stmt>     & ::= & <stmt><stmt> \\
               & \ou & <int\_var> = <int> \tok{NEWLINE} \\
    <int\_var> & ::= & int1 \ou int2 \ou int3 \\
    <int\_op>  & ::= & + \ou - \ou * \ou / \ou \% \\
    <int>      & ::= & <int\_var> \\
               & \ou & <int> <int\_op> <int> \\
  \end{tabular}
  \end{examples}
  \caption{Exemplo de uma derivação de programa GE. Adaptado de \cite{survey}.}
  \end{figure}
\end{frame}
.
% Não precisa de pacotes extras: só do \usepackage[T1]{fontenc} (já presente
% no main), para que < e > saiam corretos em fonte monoespaçada.

% Token especial (NEWLINE, INDENT) em negrito sem serifa
\providecommand{\tok}[1]{\mbox{\textsf{\textbf{#1}}}}
% Barra que separa as alternativas da BNF
\providecommand{\ou}{$\mid$\ }

\begin{frame}{\textit{Grammar-guided} GP}
  \begin{figure}
  \centering
  \begin{examples}
  \small
  \ttfamily
  \setlength{\tabcolsep}{4pt}
  \renewcommand{\arraystretch}{1.35}
  \begin{tabular}{@{}l l l@{}}
    &n = &[13, 17, 7, 24, 21, 20, 18, 28, 4, 29] \\
    <program>  & ::= & def sum(int1, int2): \tok{NEWLINE} \tok{INDENT} \\
               &     & int3 = int() \tok{NEWLINE} \\
               &     & <stmt> return <int> \\
    <stmt>     & ::= & <stmt><stmt> \\
               & \ou & <int\_var> = <int> \tok{NEWLINE} \\
    <int\_var> & ::= & int1 \ou int2 \ou int3 \\
    <int\_op>  & ::= & + \ou - \ou * \ou / \ou \% \\
    <int>      & ::= & <int\_var> \\
               & \ou & <int> <int\_op> <int> \\
  \end{tabular}
  \end{examples}
  \caption{Exemplo de uma derivação de programa GE. Adaptado de \cite{survey}.}
  \end{figure}
\end{frame}
